\section*{Цель работы}
Практическое ознакомление с простейшими схемами включения биполярных транзисторов (с общим эмиттером, общим коллектором, общей базой), а также с двумя относительно более сложными схемами (каскодной и Дарлингтона).

\section*{Теоретические сведения}

\begin{enumerate}
    \item \textit{Схема с общим эмиттером (ОЭ).} 
    Входным электродом является база, входной сигнал приложен к переходу эмиттер--база. Выходным электродом является коллектор. Эмиттер является общим электродом для входного и выходного сигналов. Схема с ОЭ усиливает сигнал как по напряжению, так и по току. Сигнал на выходе приобретает фазовый сдвиг $180^\circ$.
    
    \begin{figure}[H]
        \centering
        \includegraphics[width=0.45\linewidth]{figs/scheme_oe.pdf}
        \caption{Схема включения транзистора с общим эмиттером}
        \label{fig:scheme_oe}
    \end{figure}

    \item \textit{Схема с общим коллектором (ОК).}
    Входным электродом является база, выходным --- эмиттер. Коллектор накоротко соединен с источником питания и по переменному току заземлен, являясь общим электродом. Коэффициент передачи по напряжению $K_U < 1$, по току $K_I \gg 1$. Фазовый сдвиг $\Delta\varphi = 0$. Схема практически повторяет входной сигнал по напряжению, поэтому называется эмиттерным повторителем. Отличается высоким входным и низким выходным сопротивлением.
    
    \begin{figure}[H]
        \centering
        \includegraphics[width=0.45\linewidth]{figs/scheme_oc.pdf}
        \caption{Схема включения транзистора с общим коллектором}
        \label{fig:scheme_oc}
    \end{figure}

    \item \textit{Схема с общей базой (ОБ).}
    Входным электродом является эмиттер, выходным --- коллектор. База по переменному сигналу заземлена через конденсатор и является общим электродом. Схема имеет малое входное сопротивление (десятки Ом) и высокое выходное. Коэффициент усиления по напряжению $K_U \gg 1$, по току $K_I \approx 1$. Фазу сигнала схема не меняет.
    
    \begin{figure}[H]
        \centering
        \includegraphics[width=0.45\linewidth]{figs/scheme_ob.pdf}
        \caption{Схема включения транзистора с общей базой}
        \label{fig:scheme_ob}
    \end{figure}

    \item \textit{Сложные схемы.}
    Каскодная схема (ОБ + ОЭ) обладает малой выходной емкостью и стабильно работает на высоких частотах (\cref{fig:scheme_cascode}). Схема Дарлингтона позволяет получить сверхбольшой коэффициент усиления по току (\cref{fig:scheme_darlington}).
    
    \begin{figure}[H]
        \centering
        \begin{minipage}{0.45\linewidth}
            \centering
            \includegraphics[width=\linewidth]{figs/scheme_cascode.pdf}
            \caption{Каскодная схема}
            \label{fig:scheme_cascode}
        \end{minipage}
        \hfill
        \begin{minipage}{0.45\linewidth}
            \centering
            \includegraphics[width=\linewidth]{figs/scheme_darlington.pdf}
            \caption{Схема Дарлингтона}
            \label{fig:scheme_darlington}
        \end{minipage}
    \end{figure}
\end{enumerate}

\section*{Оборудование}
Лабораторный макет, генератор гармонических сигналов, магазин сопротивлений, два вольтметра переменного напряжения, осциллограф.

\section*{Порядок выполнения работы}

\subsection*{Подготовка к измерениям}
\begin{enumerate}
    \item Подайте напряжения питания $+15$~В и $-15$~В на лабораторный макет.
    \item Соберите измерительную цепь: выход генератора подключите ко входу исследуемой схемы. Параллельно входу подключите первый вольтметр. К выходу схемы подключите второй вольтметр и первый канал осциллографа. Второй канал осциллографа подключите ко входу схемы для контроля сдвига фаз.
\end{enumerate}

\subsection*{Исследование схемы с общим эмиттером (ОЭ)}
\begin{enumerate}
    \item Установите переключатель на лицевой панели макета в положение, соответствующее схеме ОЭ.
    \item \textit{Снятие амплитудной характеристики (АХ).}
    Установите на генераторе частоту $f$ в пределах $4 \dots 6$~кГц. Плавно увеличивайте входное напряжение $U_{\text{вх}}$ от $5$~мВ. Фиксируйте значения $U_{\text{вх}}$ и соответствующие им значения $U_{\text{вых}}$ в \cref{tab:ax_data}. Одновременно наблюдайте за формой выходного сигнала на осциллографе. Прекратите увеличение $U_{\text{вх}}$, как только на осциллограмме станут заметны нелинейные искажения (ограничение синусоиды). Зафиксируйте максимальное напряжение линейного участка $U_{\text{лин}\max}$.
    \item \textit{Снятие амплитудно-частотной характеристики (АЧХ).}
    Установите значение $U_{\text{вх}}$ на уровне, гарантированно не вызывающем искажений (на линейном участке АХ). Поддерживая $U_{\text{вх}} = \text{const}$, изменяйте частоту генератора $f$ в диапазоне от $20$~Гц до $2$~МГц. Измеряйте $U_{\text{вых}}$ и заносите данные в \cref{tab:achx_data}.
    \item \textit{Измерение фазового сдвига.}
    Установите частоту $f$ в середине полосы пропускания (где $U_{\text{вых}}$ максимально). По осциллограммам входного и выходного сигналов определите сдвиг фаз $\Delta\varphi$. Занесите результат в \cref{tab:params_data}.
    \item \textit{Измерение входного сопротивления $R_{\text{вх}}$.}
    Установите частоту в середине полосы пропускания и $U_{\text{вх}}$ без искажений. Включите магазин сопротивлений последовательно между генератором и входом схемы. Увеличивайте сопротивление магазина до тех пор, пока напряжение непосредственно на входе транзистора не упадет ровно в $2$ раза по сравнению с напряжением генератора. В этом случае $R_{\text{вх}}$ равно установленному на магазине сопротивлению. Данные занесите в \cref{tab:params_data}.
    \item \textit{Измерение выходного сопротивления $R_{\text{вых}}$.}
    Измерьте напряжение на выходе схемы без нагрузки $U_{\text{вых.хх}}$. Подключите магазин сопротивлений параллельно выходу схемы. Изменяйте его сопротивление до тех пор, пока выходное напряжение не упадет ровно в $2$ раза ($U_{\text{вых}} = U_{\text{вых.хх}} / 2$). Значение сопротивления магазина будет равно $R_{\text{вых}}$. Данные занесите в \cref{tab:params_data}.
\end{enumerate}

\subsection*{Исследование схем с общей базой (ОБ) и общим коллектором (ОК)}
\begin{enumerate}
    \item Переключите макет на схему ОБ. Повторите весь комплекс измерений (АХ, АЧХ, фаза, $R_{\text{вх}}$, $R_{\text{вых}}$), описанный для схемы ОЭ.
    \item Переключите макет на схему ОК. Повторите измерения. \textit{Внимание:} при снятии АХ для схемы ОК входное напряжение $U_{\text{вх}}$ может достигать единиц вольт (не превышать $5$~В).
\end{enumerate}

\section*{Указания к обработке результатов}
\begin{enumerate}
    \item По данным \cref{tab:ax_data} постройте графики амплитудных характеристик $U_{\text{вых}} = f(U_{\text{вх}})$ для трех схем на одной координатной плоскости. Определите значения $U_{\text{лин}\max}$.
    \item По данным \cref{tab:achx_data} рассчитайте коэффициенты усиления $K(f) = U_{\text{вых}} / U_{\text{вх}}$. Постройте графики АЧХ.
    \item По графикам АЧХ определите нижнюю и верхнюю граничные частоты $f_{\text{гр}}$ для каждой схемы из условия $K(f_{\text{гр}}) \approx 0.7 \max K(f)$.
    \item Сравните полученные экспериментальные параметры с теоретическими представлениями о свойствах схем ОЭ, ОБ и ОК.
\end{enumerate}